\subsection{Definition of semi-norms}

DEFINITION 1. --- Let E be a left vector space over K. A mapping p of E in \(R_{+} = \left[0, +\infty\right]\) , is called a semi-norm on E if it satisfies the following axioms:

\((\mathrm{SN}_{\mathrm{I}})\) If \(x \in \mathrm{E}\) and \(\lambda \in \mathrm{K}\) then \(p(\lambda x) = |\lambda| p(x)\) .

\((\mathrm{SN}_{\mathrm{II}})\) If \(x \in \mathrm{E}\) and \(y \in \mathrm{E}\) then \(p(x + y) \leqslant p(x) + p(y)\) .

Since \(p(x) \leqslant p(y) + p(x - y)\) and \(p(y) \leqslant p(x) + p(y - x)\) , from \(p(y - x) = p(x - y)\) , we deduce

\[
\left| p (x) - p (y) \right| \leqslant p (x - y).\tag{1}
\]

Examples. --- 1) A norm on E is a semi-norm p such that the relation \(p(x) = 0\) implies that x = 0 (I, p.~3).

\begin{enumerate}
\def\labelenumi{\arabic{enumi})}
\setcounter{enumi}{1}
\item
  For every linear form \(f\) on \(\mathbf{E}\) , the function \(x \mapsto |f(x)|\) is a semi-norm on \(\mathbf{E}\) .
\item
  If \(p_i (1 \leqslant i \leqslant n)\) is a finite set of semi-norms on \(E\) , then clearly \(p'(x) = \sup_{1 \leqslant i \leqslant n} p_i(x)\) .
\end{enumerate}

and \(p''(x) = \sum_{i=1}^{n} \alpha_i p_i(x)\) (where the \(\alpha_i\) are \(\geqslant 0\) ) are both semi-norms on E.

A mapping \(p\) of \(\mathbf{E}\) in \(\mathbf{R}_{+}\) is called an ultra-semi-norm if it satisfies \((\mathrm{SN}_{\mathrm{I}})\) and the following axiom:

\((\mathrm{SN}_{\mathrm{II}}^{\prime})\) If \(x \in \mathrm{E}\) and \(y \in \mathrm{E}\) , then \(p(x + y) \leqslant \sup(p(x), p(y))\) .

Clearly an ultra-semi-norm is a semi-norm.

To say that the absolute value on K is ultrametric (CA, VI, § 6.2) means that it is an ultra-semi-norm on the left vector space \(K_{s}\) , which is not identically zero.

PROPOSITION 1. --- Let E be a left topological vector space over K and let p be a semi-norm on E. The following conditions are equivalent :

\begin{enumerate}
\def\labelenumi{\alph{enumi})}
\item
  \(p\) is continuous in E.
\item
  \(p\) is continuous at the point 0.
\item
  \(p\) is uniformly continuous.
\item
  For each real number \(\alpha > 0\) , the set \(\mathbf{W}(p, \alpha)\) , of those \(x \in \mathbf{E}\) for which \(p(x) < \alpha\) , is open in \(\mathbf{E}\) .
\item
  There exists a real number \(\alpha > 0\) , such that \(\mathrm{W}(p, \alpha)\) is a neighbourhood of 0 in E.
\item
  For every real number \(\alpha > 0\) , the set \(\mathrm{V}(p, \alpha)\) , of those \(x \in \mathrm{E}\) for which \(p(x) \leqslant \alpha\) , is a neighbourhood of 0 in \(\mathrm{E}\) .
\end{enumerate}

In fact, the implications \(c) \Rightarrow a) \Rightarrow b) \Rightarrow d) \Rightarrow e) \Rightarrow f) \Rightarrow c)\) follow immediately from \((\mathrm{SN}_{\mathrm{I}})\) and inequality (1).

COROLLARY. --- If p is a continuous semi-norm on E and q is a semi-norm such that \(q \leqslant p\) , then q is continuous in E.

When \(p\) is an ultra-semi-norm on \(\mathbf{E}\) , then the sets \(\mathrm{W}(p, \alpha)\) and \(\mathrm{V}(p, \alpha)\) are both open and closed. For, we have seen that \(\mathrm{W}(p, \alpha)\) is open; on the other hand if \(z\) is a cluster point of \(\mathrm{W}(p, \alpha)\) , then there exists \(y \in \mathrm{W}(p, \alpha)\) such that \(p(y - z) < \alpha\) , and from \((\mathrm{SN}_{\mathrm{II}}')\) we have \(p(z) < \alpha\) , thus \(\mathrm{W}(p, \alpha)\) is closed. Also, \(\mathrm{V}(p, \alpha)\) is closed since \(p\) is continuous; further if \(p(x) \leqslant \alpha\) and \(p(y) \leqslant \alpha\) , then \(p(x + y) \leqslant \alpha\) by \((\mathrm{SN}_{\mathrm{II}}')\) , and this shows that \(\mathrm{V}(p, \alpha)\) is open.

